\annee{340}

\begin{notice}{1}{340}{ICUR n.s. I, 1419 · EDCS-39000407 · EDB31183}
\label{not:1}
\nf{Datation}{340 — directe : nom consulaire lu}
\nf{Genre}{Épitaphe ; ensembles A, B, C ; 92 lettres transmises ; 36 % restituées.}
\nf{Références}{ICUR I 1419 = ILCV 4647.}
\nf{Provenance}{Rattachement selon l’EDB : \emph{Incertae originis} (le volume I des ICVR réunit les inscriptions dont la provenance est perdue). Conservation selon l’EDB : \emph{Musei Vaticani: Lapidario Cristiano}.}
\nf{Support}{Plaque de marbre. 35 × 65 cm (h. × l.). Inscription gravée. Reproductions signalées par l’EDB : dessin, ICVR I (de Rossi) (droits réservés ; non reproduites). Fiche EDB : Filippo Piazzolla, première version 2015.}
\ntext{\lnum{1}[- - - coniu]gi dulcissim(a)e et fideli\\
\lnum{2}[quae vixit] annis XXVIII m(ensibus) II qu(a)e\\
\lnum{3}[deposita est - - - N]onas Aug(ustas) A⸢c⸣ind⸢y⸣no\\
\lnum{4}[et Procu]lo conss(ulibus)}
\nf{Témoins}{Export EDCS de 2020 : \emph{[3 coniu]gi dulcissim(a)e et fideli / [quae vixit] annis XXVIII m(ensibus) II qu(a)e / [deposita est 3 N]onas Aug(ustas) Aindno / [et Procu]lo conss(ulibus)}. EDB (EDB31183) : texte concordant aux conventions de notation près (indice de similarité 0,925).}
\nf{Apparat}{⸢c⸣ : Q gravé ; ⸢y⸣ : I gravé}
\nf{Traduction}{[…] à son épouse très douce et fidèle, qui vécut 28 ans 2 mois, qui fut déposée […] (jour des) nones d’août, sous le consulat d’Acindynus et de Proculus.}
\nf{Langue}{L’éditeur signale 2 graphies corrigées : Q gravé pour c ; I gravé pour y (graphie de y).}
\nf{Chronologie}{Noms consulaires : Acindynus (lu) et Proculus (restitué). Dates du jour : déposition : … Nonas Augustas (partiellement restitué). Âge déclaré : 28 a. 2 m. (partiellement restitué).}
\nf{Bibliographie}{ICUR I 1419 ; ILCV 4647. Bibliographie complémentaire selon l’EDB : ILCV, 4647 ; Ferrua 1979b, 36, n. 1419 ; Lega 2000, 105.}
\nf{Statut}{Lecture secondaire (EDCS) concordante avec l’EDB : lecture stabilisée entre deux bases, non collationnée sur l’ICVR ni sur le monument. L’EDCS signale une photographie (non consultée).}
\end{notice}

\begin{notice}{2}{340}{ICUR n.s. III, 8417 · EDCS-35802005 · EDB25692}
\label{not:2}
\nf{Datation}{340 — directe : lecture partielle d’un nom consulaire}
\nf{Genre}{Épitaphe ; ensembles A, B, C ; 46 lettres transmises ; 41 % restituées.}
\nf{Références}{ICUR III 8417.}
\nf{Provenance}{Rattachement selon l’EDB : \emph{Incertum locum in coem. Domitillae}.}
\nf{Support}{Plaque de marbre. Inscription gravée. Reproductions signalées par l’EDB : photographie PCAS ×2 (droits réservés ; non reproduites). Fiche EDB : Carlo Carletti / Antonio Enrico Felle, première version 2003.}
\ntext{\lnum{1}q]ui vixit [- - -]\\
\lnum{2}[- - -] pater fec[it - - -]\\
\lnum{3}[Acin]dino et Procu[lo conss(ulibus)}
\nf{Témoins}{EDB (EDB25692) : texte concordant aux conventions de notation près (indice de similarité 0,981).}
\nf{Traduction}{[…] qui vécut […] ; […] son père a fait (ce monument) […] ; sous le consulat d’Acindynus et de Proculus.}
\nf{Langue}{1 abréviation développée par l’éditeur (dont conss(ulibus)).}
\nf{Chronologie}{Noms consulaires : Acindynus (lu en partie) et Proculus (lu en partie).}
\nf{Bibliographie}{ICUR III 8417.}
\nf{Statut}{Lecture secondaire (EDCS) concordante avec l’EDB : lecture stabilisée entre deux bases, non collationnée sur l’ICVR ni sur le monument. L’EDCS signale une photographie (non consultée).}
\end{notice}

\begin{notice}{3}{340}{ICUR n.s. IV, 11089 · EDCS-37200371 · EDB39548}
\label{not:3}
\nf{Datation}{340 — directe : lecture partielle d’un nom consulaire}
\nf{Genre}{Épitaphe ; ensembles A, B, C ; 58 lettres transmises ; 64 % restituées.}
\nf{Références}{ICUR IV 11089.}
\nf{Provenance}{Rattachement selon l’EDB : \emph{Coem. Callisti area subdialis et incerti loci}. Conservation selon l’EDB : \emph{Coem. Callisti area subdialis et incerti loci}.}
\nf{Support}{Plaque de marbre. 18 × 22 × 4,5 cm (h. × l. × ép.). Inscription gravée. Fiche EDB : Marida Pierno, première version 2015.}
\ntext{\lnum{1}ben]e meren[ti in pace - - -]\\
\lnum{2}[- - - perii]t pri(die) Idu[s - - -]\\
\lnum{3}[- - - Acin]dino e[t Proculo conss(ulibus)}
\nf{Témoins}{EDB (EDB39548 ; écarts à examiner) : \emph{------ / [ben]emeren[ti in pace] / [peri]t pri(die) idu[s ---] / [---]dino e[t ---] / ------}.}
\nf{Traduction}{[…] bien méritant, en paix […] ; […] mourut la veille des ides […] ; sous le consulat d’Acindynus et de Proculus.}
\nf{Langue}{2 abréviations développées par l’éditeur (dont pri(die), conss(ulibus)).}
\nf{Chronologie}{Noms consulaires : Acindynus (lu en partie) et Proculus (restitué).}
\nf{Commentaire}{La formule in pace est attestée seulement par restitution.}
\nf{Bibliographie}{ICUR IV 11089.}
\nf{Statut}{Lecture secondaire (EDCS) ; écarts avec l’EDB à arbitrer sur l’ICVR. Plus de la moitié des lettres sont restituées.}
\end{notice}

\annee{341}

\begin{notice}{4}{341}{ICUR n.s. I, 1420 · EDCS-34501143 · EDB31184}
\label{not:4}
\nf{Datation}{341 — directe : nom consulaire lu}
\nf{Genre}{Épitaphe ; ensembles A, B, C ; 189 lettres transmises.}
\nf{Références}{ICUR I 1420 = ILCV 2816 (em) = ICVaticano p 299 = ICaRoma 89 = AE 1997, 166.}
\nf{Provenance}{Rattachement selon l’EDB : \emph{Incertae originis} (le volume I des ICVR réunit les inscriptions dont la provenance est perdue). Conservation selon l’EDB : \emph{Musei Vaticani: Lapidario Cristiano}.}
\nf{Support}{Plaque de marbre. 49 × 186 cm (h. × l.). Inscription gravée. Signes décrits par l’EDB : chrisme. Reproductions signalées par l’EDB : dessin, ICVR I (de Rossi) et reproduction, Perret 1851 (droits réservés ; non reproduites). Fiche EDB : Filippo Piazzolla, première version 2015.}
\ntext{\lnum{1}Depossio Iuniani pri(die) Idus Apriles Marcellino et Probino conss(ulibus)\\
\lnum{2}qui ⸢v⸣ixit annis XL in pace dec⸢e⸣ssit et amator paupe\{ro\}rum ⸢v⸣ixit\\
\lnum{3}cum ⸢v⸣(i)rg(i)nia annis XV bene merenti ⸢v⸣irginia sua Bictora\\
\lnum{4}bene merenti fecit amatrix paupe\{ro\}rum et operaria}
\nf{Témoins}{Export EDCS de 2020 : \emph{Depossio Iuniani pri(die) Idus Apriles Marcellino et Probino conss(ulibus) / qui ixit annis XL in pace decssit et amator paupe\{ro\}rum ixit / cum (i)rg(i)nia annis XV bene merenti irginia sua Bictora / bene merenti fecit amatrix paupe\{ro\}rum et operaria}. EDB (EDB31184) : même texte ; les écarts relevés dans la liste de résultats tiennent à la notation (relecture de la fiche détaillée). Relecture de la fiche détaillée de l’EDB : Même texte. L'EDB garde la graphie de la pierre (bixit, birginia), l'EDCS la corrige (<v=B>). L'EDB décrit un chrisme final. Supports bibliographiques : Di Stefano Manzella 1997 (avec photo), Carletti 1986a n. 89.}
\nf{Apparat}{⸢v⸣ : B gravé ; ⸢e⸣ : I gravé ; ⸢v⸣ : B gravé ; ⸢v⸣ : B gravé ; ⸢v⸣ : B gravé ; \{ro\} : lettres superflues ; \{ro\} : lettres superflues}
\nf{Traduction}{Déposition de Iunianus, la veille des ides d’avril, sous le consulat de Marcellinus et de Probinus ; il vécut 40 ans ; il est mort en paix ; ami des pauvres, il vécut avec sa vierge (= épouse épousée vierge) 15 ans ; à lui, bien méritant, sa vierge Victora a fait (ce monument), bien méritant ; amie des pauvres et travailleuse.}
\nf{Langue}{L’éditeur signale 5 graphies corrigées : B gravé pour v (bêtacisme) ×4 ; I gravé pour e (confusion e/i). Des lettres gravées sont jugées superflues par l’éditeur (entre accolades). 4 abréviations développées par l’éditeur (dont pri(die), conss(ulibus), (i)rg, (i)nia).}
\nf{Chronologie}{Noms consulaires : Marcellinus (lu) et Probinus (lu). Dates du jour : déposition : pridie Idus Apriles = 12 avril 341 (lu). Âge déclaré : 40 a. (lu).}
\nf{Commentaire}{La formule in pace est conservée. Commentaire développé : dossier p.~\pageref{dos:4}.}
\nf{Bibliographie}{ICUR I 1420 ; ILCV 2816 (em) ; ICVaticano p 299 ; ICaRoma 89 ; AE 1997, 166. Bibliographie complémentaire selon l’EDB : Di Stefano Manzella 1997, 299-200, 3.7.9, con foto (C. Lega) ; Lega 2000, 92 ; Carletti 1986a, 101-102, n. 89 ; ILCV, 2816.}
\nf{Statut}{Lecture secondaire (EDCS) concordante avec l’EDB : lecture stabilisée entre deux bases, non collationnée sur l’ICVR ni sur le monument. L’EDCS signale une photographie (non consultée).}
\end{notice}

\annee{342}

\begin{notice}{5}{342}{ICUR n.s. I, 1251 · EDCS-39101961 · EDB30459}
\label{not:5}
\nf{Datation}{342 — directe : lecture partielle d’un nom consulaire}
\nf{Genre}{Épitaphe ; ensembles A, B, C ; 98 lettres transmises ; 55 % restituées.}
\nf{Références}{ICUR I 1251 = Louvre 966.}
\nf{Provenance}{Rattachement selon l’EDB : \emph{Incertae originis} (le volume I des ICVR réunit les inscriptions dont la provenance est perdue). Conservation selon l’EDB : \emph{Musei Capitolini}.}
\nf{Support}{Plaque de marbre. 26 × 30 cm (h. × l.). Inscription gravée. Fiche EDB : Filippo Piazzolla, première version 2014.}
\ntext{\lnum{1}vixit ann(os) - - -]XVIII m(enses) [- - -]\\
\lnum{2}[- - -] Non(as) Sept[embr(es) - - -]\\
\lnum{3}[dd(ominis) nn(ostris) Co]nstanti[o III]\\
\lnum{4}[et Const]ante II Augg(ustis) [cons(ulibus)]\\
\lnum{5}[The]odotus [posuit}
\nf{Témoins}{EDB (EDB30459 ; écarts à examiner) : \emph{------ / [--- vixit annis] XVIII m[ensibus ---] / [--- dep(osit---) ---] ṇon(as) sept[embres ---] / [dd(ominis) nn(ostris) Co]nstanti[o III] / [et Const]ante II Augg̣[ustibus cons(ulibu)s] / [--- The]odotus [--- posuit] / ------ [texte tronqué dans la liste de résultats]}.}
\nf{Traduction}{[…] vécut […]18 ans, […] mois ; […] nones de septembre […] ; sous le consulat de nos seigneurs Constance (Auguste) pour la 3e fois et Constant (Auguste) pour la 2e fois ; Theodotus a posé (ce monument).}
\nf{Langue}{8 abréviations développées par l’éditeur (dont ann(os), m(enses), Non(as), Septembr(es), dd(ominis), nn(ostris)).}
\nf{Chronologie}{Noms consulaires : Constantius (lu en partie) et Constans (lu en partie). Dates du jour : événement indéterminé : … Nonas Septembres (partiellement restitué).}
\nf{Commentaire}{Vocabulaire religieux : anthroponyme (niveau 6) — theodotus (anthroponyme).}
\nf{Bibliographie}{ICUR I 1251 ; Louvre 966. Bibliographie complémentaire selon l’EDB : Molisani 1973, 110 ; Ferrua 1979b, 30, n. 1251.}
\nf{Statut}{Lecture secondaire (EDCS) ; écarts avec l’EDB à arbitrer sur l’ICVR. Plus de la moitié des lettres sont restituées. L’EDCS signale une photographie (non consultée).}
\end{notice}

\begin{notice}{6}{342}{ICUR n.s. IV, 9556 · EDCS-36101085 · EDB20872}
\label{not:6}
\nf{Datation}{342 — directe : lecture partielle d’un nom consulaire}
\nf{Genre}{Épitaphe ; ensembles A, B, C ; 152 lettres transmises ; 66 % restituées.}
\nf{Références}{ICUR IV 9556 = ILCV +2846.}
\nf{Provenance}{Rattachement selon l’EDB : \emph{Coem. Callisti pars inferior}.}
\nf{Support}{Plaque de marbre. Inscription gravée. Reproductions signalées par l’EDB : photographie PCAS (droits réservés ; non reproduites). Fiche EDB : Antonio Enrico Felle, première version 2003.}
\ntext{\lnum{1}[Aur(elius)] Paulinus Aur(eliae) Ta[eo]dore [filiae dul]\\
\lnum{2}[ciss]im(a)e fecit qu(a)e pr[aec]essit in [pace dd(ominis) nn(ostris)]\\
\lnum{3}[Con]statio Aug(usto) III et [Constante Aug(usto) II conss(ulibus)\\
\lnum{4}[- - - D]ec(embres) et ⸢v⸣ixit an(nos) LI[- - -]}
\nf{Témoins}{Export EDCS de 2020 : \emph{[Aur(elius)] Paulinus Aur(eliae) Ta[eo]dore [filiae dul]/[ciss]im(a)e fecit qu(a)e pr[aec]essit in [pace dd(ominis) nn(ostris)] / [Con]statio Aug(usto) III et [Constante Aug(usto) II conss(ulibus) / [3 D]ec(embres) et ixit an(nos) LI[3]}. EDB (EDB20872) : texte concordant aux conventions de notation près (indice de similarité 0,98).}
\nf{Apparat}{⸢v⸣ : B gravé}
\nf{Traduction}{Aurelius Paulinus à Aurelia Theodora, sa fille très douce, a fait (ce monument) ; elle (l’)a précédé dans la paix, sous le consulat de nos seigneurs Constance Auguste pour la 3e fois et Constant Auguste pour la 2e fois, […] décembre ; elle vécut 51 (?) ans […].}
\nf{Langue}{L’éditeur signale une graphie corrigée : B gravé pour v (bêtacisme).}
\nf{Chronologie}{Noms consulaires : Constantius (lu en partie) et Constans (restitué).}
\nf{Commentaire}{La formule in pace est attestée seulement par restitution.}
\nf{Bibliographie}{ICUR IV 9556 ; ILCV +2846.}
\nf{Statut}{Lecture secondaire (EDCS) concordante avec l’EDB : lecture stabilisée entre deux bases, non collationnée sur l’ICVR ni sur le monument. Plus de la moitié des lettres sont restituées. L’EDCS signale une photographie (non consultée).}
\end{notice}

\annee{343}

\begin{notice}{7}{343}{ICUR n.s. I, 264 · EDCS-38600270 · EDB27782}
\label{not:7}
\nf{Datation}{343 — directe : nom consulaire lu}
\nf{Genre}{Épitaphe ; ensembles A, B, C ; 86 lettres transmises ; 12 % restituées.}
\nf{Références}{ICUR I 264 = ILCV 4393.}
\nf{Provenance}{Rattachement selon l’EDB : \emph{Incertae originis} (le volume I des ICVR réunit les inscriptions dont la provenance est perdue).}
\nf{Support}{Support non précisé par l’EDB. Reproductions signalées par l’EDB : dessin, ICVR I (de Rossi) (droits réservés ; non reproduites). Fiche EDB : Filippo Piazzolla, première version 2014.}
\ntext{\lnum{1}Placi⸢d⸣o ⸢Fl⸣(avio) Romulo consulibus III Nona[s Iulias die]\\
\lnum{2}Martis decessit Leo annorum XXX PO[- - -]\\
\lnum{3}SV vixit ann(os) nn(u)mm(ero ?) X}
\nf{Témoins}{Export EDCS de 2020 : \emph{Placio (avio) Romulo consulibus III Nona[s Iulias die] / Martis decessit Leo annorum XXX PO[3] / SV vixit ann(os) nn(u)mm(ero ?) X}. EDB (EDB27782) : texte concordant aux conventions de notation près (indice de similarité 0,948).}
\nf{Apparat}{⸢d⸣ : T gravé ; ⸢Fl⸣ : IT gravé}
\nf{Traduction}{Sous le consulat de Placidus et de Flavius Romulus, le 3e jour avant les nones de juillet, un mardi, mourut Leo, âgé de 30 ans […] ; […] vécut en tout 10 (?) ans [texte corrompu].}
\nf{Langue}{L’éditeur signale 2 graphies corrigées : T gravé pour d ; IT gravé pour Fl.}
\nf{Chronologie}{Noms consulaires : Placidus (lu) et Romulus (lu). Dates du jour : événement indéterminé : III Nonas Iulias = 5 juillet 343 (partiellement restitué). Âge déclaré : 30 a. (lu).}
\nf{Bibliographie}{ICUR I 264 ; ILCV 4393. Bibliographie complémentaire selon l’EDB : ILCV, 4393.}
\nf{Statut}{Lecture secondaire (EDCS) concordante avec l’EDB : lecture stabilisée entre deux bases, non collationnée sur l’ICVR ni sur le monument. L’EDCS signale une photographie (non consultée).}
\end{notice}

\begin{notice}{8}{343}{ICUR n.s. I, 2080 · EDCS-25402137 · EDB15588}
\label{not:8}
\nf{Datation}{343 — directe : nom consulaire lu}
\nf{Genre}{Épitaphe ; ensembles A, B, C ; 37 lettres transmises.}
\nf{Références}{ICUR I 2080.}
\nf{Provenance}{Rattachement selon l’EDB : \emph{Incertae originis} (le volume I des ICVR réunit les inscriptions dont la provenance est perdue). Conservation selon l’EDB : \emph{Musei Vaticani: Lapidario Cristiano}.}
\nf{Support}{Sarcophage de marbre. 15 × 20 cm (h. × l.). Inscription gravée. Fiche EDB : Antonio Enrico Felle, première version 2003.}
\ntext{\lnum{1}κατάθεσις\\
\lnum{2}Ὀκταβίλλης\\
\lnum{3}πρὸ θʹ καλ(ανδῶν) σε(π)τ(εμβρίων)\\
\lnum{4}(ὑπατείᾳ) Ῥωμουλι}
\nf{Témoins}{Export EDCS de 2020 : \emph{texte grec absent}. EDB (EDB15588 ; écarts à examiner) : \emph{κατάθεσις / Ὀκταβίλλης / πρὸ θ´ καλ(ανδῶν) σε(π)τ(εμβρίων) / ⟨ὑπατείᾳ⟩ Ῥωμουλι [texte tronqué dans la liste de résultats]}.}
\nf{Traduction}{(grec) Déposition d’Octavilla, le 9e jour avant les calendes de septembre, sous le consulat de Romulus.}
\nf{Langue}{Le texte est rédigé en grec, en tout ou en partie ; l’export de 2020 ne le transmettait pas. 4 abréviations développées par l’éditeur (dont (ανδῶν), (π), (εμβρίων), (ὑπατείᾳ)).}
\nf{Chronologie}{Noms consulaires : Romulus (lu).}
\nf{Bibliographie}{ICUR I 2080. Bibliographie complémentaire selon l’EDB : Deichmann 1967, 70, no. 84 (with photo at taf. 24) ; CLRE, 633.}
\nf{Statut}{Lecture secondaire (EDCS) ; écarts avec l’EDB à arbitrer sur l’ICVR.}
\end{notice}

\begin{notice}{9}{343}{ICUR n.s. III, 8720 · EDCS-35900965 · EDB18778}
\label{not:9}
\nf{Datation}{343 — directe : nom consulaire lu}
\nf{Genre}{Épitaphe ; ensembles A, B, C ; 68 lettres transmises.}
\nf{Références}{ICUR III 8720 = ILCV 2571 (em) = ILCV 2819 (em).}
\nf{Provenance}{Rattachement selon l’EDB : \emph{Tituli qui in coem. Callisti reperti traduntur}.}
\nf{Support}{Plaque de marbre. Inscription gravée. Signes décrits par l’EDB : chrisme. Reproductions signalées par l’EDB : dessin, ICVR I (de Rossi) (droits réservés ; non reproduites). Fiche EDB : Carlo Carletti, première version 2003.}
\ntext{\lnum{1}Placido et Romulo (!) XVII ⸢K⸣(alendas) Dec(embres)\\
\lnum{2}Ianuaria in pace quae\\
\lnum{3}vixit annos XVIIII m(enses)\\
\lnum{4}VIII d(ies) XIIII}
\nf{Témoins}{Export EDCS de 2020 : \emph{Placido et Romulo (!) XVII (alendas) Dec(embres) / Ianuaria in pace quae / vixit annos XVIIII m(enses) / VIII d(ies) XIIII}. EDB (EDB18778) : texte concordant aux conventions de notation près (indice de similarité 0,988).}
\nf{Apparat}{⸢K⸣ : C gravé ; (!) : forme signalée par l’éditeur}
\nf{Traduction}{Sous (le consulat de) Placidus et Romulus, le 17e jour avant les calendes de décembre, Ianuaria, en paix, qui vécut 19 ans 8 mois 14 jours.}
\nf{Langue}{L’éditeur signale une graphie corrigée : C gravé pour K. L’éditeur signale au moins une forme anormale (sic). 4 abréviations développées par l’éditeur (dont (alendas), Dec(embres), m(enses), d(ies)).}
\nf{Chronologie}{Noms consulaires : Placidus (lu) et Romulus (lu). Dates du jour : événement indéterminé : XVII Kalendas Decembres = 15 novembre 343 (lu). Âge déclaré : 19 a. 8 m. 14 j. (lu).}
\nf{Commentaire}{La formule in pace est conservée.}
\nf{Bibliographie}{ICUR III 8720 ; ILCV 2571 (em) ; ILCV 2819 (em).}
\nf{Statut}{Lecture secondaire (EDCS) concordante avec l’EDB : lecture stabilisée entre deux bases, non collationnée sur l’ICVR ni sur le monument. L’EDCS signale une photographie (non consultée).}
\end{notice}

\annee{344}

\begin{notice}{10}{344}{ICUR n.s. IV, 11752 · EDCS-37400847 · EDB699}
\label{not:10}
\nf{Datation}{344 — directe : lecture partielle d’un nom consulaire}
\nf{Genre}{Épitaphe ; ensembles A, B, C ; 188 lettres transmises ; 66 % restituées.}
\nf{Références}{ICUR IV 11752.}
\nf{Provenance}{Rattachement selon l’EDB : \emph{Coem. anonymum ad viam Appiam}. Conservation selon l’EDB : \emph{Coem. anonymum ad viam Appiam}.}
\nf{Support}{Plaque de marbre. Inscription gravée. Fiche EDB : Maria Cristina Grisanzio, première version 2003.}
\ntext{\lnum{1}[Her]mogenia [quae vixit a]nnis XXVI m[ensibus - - - diebus - - -]\\
\lnum{2}[recessit - - - Non(as) Mai(as) deposi]ta Nonis s(upra) s(criptis) u[xor Primi - - -]\\
\lnum{3}[hic requi]escit in pace [- - -] Primus co[mpari su(a)e fecit]\\
\lnum{4}[bene] merenti in pac[e Leonti]o et Salu[stio vv(iris) cc(larissimis) co(n)ss(ulibus)]}
\nf{Témoins}{EDB (EDB699 ; écarts à examiner) : \emph{[Her]mogenia [quae vixit a]nnis XXVI m[ensibus --- diebus ---] / [--- deposi]ta nonis s(upra) s(criptis) u[xor Primi ---] / [hic requi]escit in pace --- Primus co[mpari suae fecit] / [bene]merenti in pac[e Leonti]o et Salu[stio viris clarissimis co(n)s(ulibu)s] [texte tronqué dans la liste de résultats]}.}
\nf{Traduction}{Hermogenia, qui vécut 26 ans […] mois […] jours ; elle s’en alla […] (avant) les nones de mai ; déposée aux nones susdites ; épouse de Primus […] ; elle repose ici en paix ; […] Primus à sa compagne a fait (ce monument), bien méritante, en paix ; sous le consulat de Leontius et Sallustius, clarissimes.}
\nf{Langue}{9 abréviations développées par l’éditeur (dont Non(as), Mai(as), s(upra), s(criptis), su(a)e, vv(iris)).}
\nf{Chronologie}{Noms consulaires : Leontius (restitué) et Sallustius (absent). Dates du jour : décès : … Nonas Maias (restitué). Âge déclaré : 26 a. (partiellement restitué).}
\nf{Commentaire}{La formule in pace est conservée. Vocabulaire religieux : repos et sommeil de la mort — requiescit (partiellement restitué).}
\nf{Bibliographie}{ICUR IV 11752.}
\nf{Statut}{Lecture secondaire (EDCS) ; écarts avec l’EDB à arbitrer sur l’ICVR. Plus de la moitié des lettres sont restituées.}
\end{notice}

\begin{notice}{38}{352}{ICUR n.s. II, 4241 · EDCS-33900148 · EDB17657}
\label{not:38}
\nf{Datation}{352 — directe : nom consulaire lu}
\nf{Genre}{Épitaphe ; ensembles A, B, C ; 103 lettres transmises.}
\nf{Références}{ICUR II 4241 = ILCV 3252a (em).}
\nf{Provenance}{Rattachement selon l’EDB : \emph{Coem. sub divo apud basilicam S.Petri apostoli}. Conservation selon l’EDB : \emph{Musei Vaticani: Lapidario Cristiano}.}
\nf{Support}{Plaque de marbre. Inscription gravée. Signes décrits par l’EDB : chrisme. Reproductions signalées par l’EDB : dessin, ICVR I (de Rossi) (droits réservés ; non reproduites). Fiche EDB : Antonio Enrico Felle, première version 2003.}
\ntext{\lnum{1}Privata qu(a)e vixit annis LIII me(n)ses(!)\\
\lnum{2}VIII die(bu)s II deposita est IIII Nonas\\
\lnum{3}Febr(uarias) bene merenti in pace Chr(isti)\\
\lnum{4}Decentio et Paulo conss(ulibus)}
\nf{Témoins}{EDB (EDB17657) : texte concordant aux conventions de notation près (indice de similarité 0,972).}
\nf{Apparat}{(!) : forme signalée par l’éditeur}
\nf{Traduction}{Privata, qui vécut 53 ans 8 mois 2 jours ; déposée le 4e jour avant les nones de février ; bien méritante, en paix ☧ ; sous le consulat de Décence et Paulus.}
\nf{Langue}{Cinq autres développements : qu(a)e, me(n)ses, die(bu)s, Febr(uarias), conss(ulibus). Chr(isti) est conservé comme transmission éditoriale ; son statut de mot abrégé n’est pas tenu pour établi.}
\nf{Chronologie}{Noms consulaires : Decentius (lu) et Paulus (lu). Dates du jour : déposition : IIII Nonas Februarias = 2 février 352 (lu). Âge déclaré : 53 a. 8 m. 2 j. (lu).}
\nf{Commentaire}{Après \emph{in pace}, l'EDCS transmet \emph{Chr(isti)} ; l'EDB (EDB17657, notice d'A. E. Felle) transcrit \emph{in pace ((monogramma Christi))} et décrit un chrisme dans les données matérielles. L'estampage universitaire ILCV 3252A (Ohio State University) montre lui aussi un christogramme placé après IN PACE. Deux témoins indépendants de l'EDCS concordent donc : la pierre porte un monogramme, non le mot \emph{Christi} écrit en lettres. Le texte est classé parmi les signes (catégorie S) et non parmi les mentions lexicales du Christ. La question de savoir si le monogramme fonctionne ici comme abréviation (\emph{in pace Chr(isti)}) ou comme signe juxtaposé ne peut être tranchée sans photographie : elle est discutée avec les n° 140, 151 et 311 (chapitre sur Dieu et le Christ). Commentaire développé : dossier p.~\pageref{dos:38}.}
\nf{Bibliographie}{ICUR II 4241 ; ILCV 3252a (em).}
\nf{Statut}{Lecture secondaire (EDCS) concordante avec l’EDB : lecture stabilisée entre deux bases, non collationnée sur l’ICVR ni sur le monument. L’EDCS signale une photographie (non consultée).}
\end{notice}

\begin{notice}{106}{367}{EDCS-16600352 · EDB13167 · EDH HD030246}
\label{not:106}
\nf{Statut}{Doublon établi : même défunt (Zosimus), même date (13 mai 367), même formule \emph{annorum plus minus}, même coupe inhabituelle des lignes (\emph{Ma/ias}, \emph{Lupicin/o}). L'EDH (HD030246) et l'EDB (EDB13167) rattachent le texte de l'AE 1905, 71 à ICUR II 6028. La notice est fusionnée dans le n° 120 ; sa leçon \emph{c(iu)esq(it!)}, contre \emph{recessit} dans l'ICVR et l'EDB, est donnée en apparat du n° 120.}
\nf{Références}{AE 1905, 71}
\ntext{\lnum{1}Zosimus c(iu)esq(it !) in pace\\
\lnum{2}annorum plus minus\\
\lnum{3}[- - -] III Idus Ma\\
\lnum{4}[ias] Lupicin\\
\lnum{5}o et Io⸢v⸣ino cons(ulibus)}
\nf{Traduction}{Zosimus repose en paix, âgé d’environ […] ans ; […] le 3e jour avant les ides de mai, sous le consulat de Lupicinus et de Iovinus.}
\nf{Voir}{la notice 120 (p.~\pageref{not:120}).}
\end{notice}

\begin{notice}{214}{380}{ICUR n.s. I, 3194 · EDCS-35200283 · EDB33145}
\label{not:214}
\nf{Datation}{380 — formule transmise par l’EDB (lue)}
\nf{Genre}{Épitaphe ; ensembles A, B, C ; 129 lettres transmises.}
\nf{Références}{CLE 1619 = ICUR I 3194 = ILCV 1464.}
\nf{Provenance}{Rattachement selon l’EDB : \emph{Incertae originis} (le volume I des ICVR réunit les inscriptions dont la provenance est perdue). Conservation selon l’EDB : \emph{villa Albani}.}
\nf{Support}{Plaque de marbre. 69 × 128 cm (h. × l.). Inscription gravée. Reproductions signalées par l’EDB : reproduction, Perret 1851 (droits réservés ; non reproduites). Fiche EDB : Filippo Piazzolla, première version 2016.}
\ntext{\lnum{1}Hic quiescit ancilla dei qu(a)e de\\
\lnum{2}sua omnia possedit domum ista\\
\lnum{3}qu⸢a⸣m amic(a)e deflen(t) solaciumq(ue) requirunt\\
\lnum{4}pro hunc unum ora subolem qu⸢a⸣m superis\\
\lnum{5}titem re(li)quisti}
\nf{Témoins}{Export EDCS de 2020 : \emph{Hic quiescit ancilla dei qu(a)e de / sua omnia possedit domum ista / qum amic(a)e deflen(t) solaciumq(ue) requirunt / pro hunc unum ora subolem qum superis/titem re(li)quisti}. EDB (EDB33145 ; divergence forte) : \emph{hic quescit ancilla dei ⌜q⌝ue de / sua omnia possedit domum ista / quem amice defien solaciumq(ue) requirun̂t / pro hunc unum ora subolem quem superis / titem requisti eterna requiem felicita/⟨ti⟩s cuasa man(e)bis IIIIX k⌜a⌝lendas octobris / Cucurbitinus et Abu⌜nd⌝antius hic ⌜s⌝imul quiescit / dd(ominis) nn(ostris) Gratiano V et Teodosio aaugg(ustis) ((aves coronam tenentes)) [texte tronqué dans la liste de résultats]}. Traduction des lignes transmises par l'EDB seule : « … tu demeureras dans le repos éternel de la félicité (?) ; le … jour avant les calendes d'octobre. Cucurbitinus et Abundantius reposent ici ensemble. Sous le consulat de nos seigneurs Gratien, pour la cinquième fois, et Théodose, Augustes. »}
\nf{Apparat}{⸢a⸣ : E gravé ; ⸢a⸣ : E gravé}
\nf{Traduction}{Ici repose une servante de Dieu qui, de ses biens, posséda toute cette maison, que ses amies pleurent et dont elles cherchent le réconfort ; prie pour ce seul (enfant), la descendance que tu as laissée survivante.}
\nf{Langue}{L’éditeur signale 2 graphies corrigées : E gravé pour a (e pour a) ×2.}
\nf{Chronologie}{Le texte EDCS s'arrête à \emph{re(li)quisti}. La fiche EDB (EDB33145) transmet trois lignes supplémentaires : \emph{eterna requiem felicita⟨ti⟩s cuasa man(e)bis IIIIX k⌜a⌝lendas octobris / Cucurbitinus et Abu⌜nd⌝antius hic ⌜s⌝imul quiescit / dd(ominis) nn(ostris) Gratiano V et Teodosio aaugg(ustis)}. L'année 380 est donc fondée sur une formule consulaire lue. Le quantième \emph{IIIIX} est anomal. Deux autres défunts, Cucurbitinus et Abundantius, reposent dans la même tombe ; l'EDB décrit des oiseaux tenant une couronne. La plaque (69 × 128 cm) est conservée à la villa Albani.}
\nf{Commentaire}{Vocabulaire religieux : désignation de Dieu — dei (lu), ancilla dei (lu) ; intercession — ora (lu) ; repos et sommeil de la mort — quiescit (lu). Désignation divine : retenue dans le profil strict. Commentaire développé : dossier p.~\pageref{dos:214}.}
\nf{Bibliographie}{CLE 1619 ; ICUR I 3194 ; ILCV 1464. Bibliographie complémentaire selon l’EDB : ILCV, 1464.}
\nf{Statut}{Lecture secondaire (EDCS) ; divergence forte avec l’EDB : le texte ne doit pas être exploité sans contrôle de l’édition imprimée. L’EDCS signale une photographie (non consultée).}
\end{notice}

\begin{notice}{447}{401}{ICUR n.s. I, 1942 · EDCS-25402147 · EDB12500}
\label{not:447}
\nf{Datation}{401 — formule transmise par l’EDB (partiellement restituée)}
\nf{Genre}{Épitaphe ; ensembles A, B, C ; 122 lettres transmises ; 43 % restituées.}
\nf{Références}{AE 1910, 205 = ICUR I 1942 = SEG 30, 1208.}
\nf{Provenance}{Rattachement selon l’EDB : \emph{Incertae originis} (le volume I des ICVR réunit les inscriptions dont la provenance est perdue). Conservation selon l’EDB : \emph{Museo Archeologico Nazionale}.}
\nf{Support}{Plaque de marbre. 34 × 26 × 3 cm (h. × l. × ép.). Inscription gravée. Fiche EDB : Antonio Enrico Felle, première version 2003.}
\ntext{\lnum{1}ὁρκίζομε δὲ ὑμᾶς [τὸν θ(εὸ)ν τὸν πα]\\
\lnum{2}ντοκράτορα εἴ τις το[λμήσει σκυ]\\
\lnum{3}λῖσε τὰ ὀστέα αὐτοῦ [ἢ ἄλλον τινὰ]\\
\lnum{4}[θῖνε ἔσ]τε ἔνοχος τῇ [κρίσει - - -]\\
\lnum{5}[- - -]αν ἡμαίραν τῆς [ἀναστάσεως(?)]}
\nf{Témoins}{Export EDCS de 2020 : \emph{texte grec absent}. EDB (EDB12500 ; divergence forte) : \emph{[ἐνθά]δε κεῖτε Εὐ[---] / [Θεο]δοσειοπολίτ[ης ἀμέμπτως ζήσ]/[ας] μικρῷ πλέω ἔ[τη --- ἐκοιμήθη τ]/ῇ πρὸ ιε´ καλ(ανδῶν) σεπτ(εμβρίων) [ὑπατίᾳ Βινκεντί]/ου καὶ Φραβίτου τῶ[ν λαμπροτάτων] / ὁρκίζομε δὲ ὑμᾶς [διὰ τὸν θεὸν τὸν πα]/ντοκράτορα εἴ τις το]λμήσειε σκυ]/λῖσε τὰ ᾡτέα αὐτοῦ [ἢ τάφον τοῦτο]/[ν βλάψε]ιε ἔνοχος τῇ [κολάσει τὴν φο]/[βερ]ὰν ἡμαίραν τῆς [κρίσεως ἔστω] [texte tronqué dans la liste de résultats]}. Les deux témoins diffèrent aussi dans les restitutions. Début (EDB seule) : \emph{[ἐνθά]δε κεῖτε Εὐ[---] / [Θεο]δοσειοπολίτ[ης]} : le défunt est originaire de Théodosiopolis. Fin : l'EDCS restitue \emph{ἔνοχος τῇ [κρίσει - - -] / [- - -]αν ἡμαίραν τῆς [ἀναστάσεως(?)]} (« passible du jugement … au jour de la résurrection ») ; l'EDB restitue \emph{ἔνοχος τῇ [κολάσει τὴν φο]/[βερ]ὰν ἡμαίραν τῆς [κρίσεως ἔστω]} (« qu'il soit passible du châtiment au jour terrible du jugement »). Seuls \emph{ἔνοχος τῇ} et \emph{ἡμαίραν τῆς} sont lus. La mention de la résurrection, déjà conjecturale, est donc contestée par la seconde base : elle ne peut entrer dans aucun compte.}
\nf{Traduction}{(grec, texte EDCS) Je vous adjure par [Dieu le Tout]-Puissant : si quelqu'un ose troubler ses os [ou (déposer) quelqu'un d'autre], qu'il soit passible du [jugement …] au jour de la [résurrection (?)]. — Pour le texte plus complet et les restitutions différentes de l'EDB, voir la confrontation des témoins.}
\nf{Langue}{Le texte est rédigé en grec, en tout ou en partie ; l’export de 2020 ne le transmettait pas. 1 abréviation développée par l’éditeur (dont (εὸ)).}
\nf{Chronologie}{Le texte EDCS ne transmet que les cinq dernières lignes. La fiche EDB (EDB12500 = ICVR I 1942), qui suit une autre édition, en donne dix, dont la date : \emph{τῇ πρὸ ιε´ καλ(ανδῶν) σεπτ(εμβρίων) [ὑπατίᾳ Βινκεντί]ου καὶ Φραβίτου} (18 août 401). Le nom de Fravitta est lu : l'année 401 est fondée sur une formule consulaire partiellement restituée.}
\nf{Commentaire}{Vocabulaire religieux (texte EDCS) : désignation de Dieu — \emph{θ(εὸ)ν} (restitué), \emph{-ντοκράτορα} (lu en partie) ; jugement — \emph{κρίσει} (restitué) ; résurrection — \emph{ἀναστάσεως} (restitué avec doute, contesté par l'EDB). L'adjuration (\emph{ὁρκίζομε δὲ ὑμᾶς}) et l'épithète \emph{-ντοκράτορα} sont lues ; le nom de Dieu est restitué. Le texte est une formule de protection de la tombe. L'idée d'un jugement (\emph{κρίσις}) apparaît dans les deux restitutions, à deux places différentes ; la résurrection, dans une seule. Commentaire développé : dossier p.~\pageref{dos:447}.}
\nf{Bibliographie}{AE 1910, 205 ; ICUR I 1942 ; SEG 30, 1208.}
\nf{Statut}{Lecture secondaire (EDCS) ; divergence forte avec l’EDB : le texte ne doit pas être exploité sans contrôle de l’édition imprimée.}
\end{notice}